% Source: https://artofproblemsolving.com/community/c3h2125065_400_posts_lets_go
% Stable identity: aops-topic-2125065-q01
\textbf{(a)} Take $x=y=-1$. Then $\sqrt{xy}=\sqrt1=1$, but $\dfrac{x+y}{2}=-1<1$.

\textbf{(b)} Every step up to $\dfrac{(x+y)^2}{4}\ge xy$ is valid for all real $x$ and $y$. The last step is the problem. If $u\ge v\ge0$, then $\sqrt u\ge\sqrt v$; but
\[\sqrt{\dfrac{(x+y)^2}{4}}=\dfrac{|x+y|}{2},\]
not $\dfrac{x+y}{2}$. So the argument really proves $\dfrac{|x+y|}{2}\ge\sqrt{xy}$ whenever $xy\ge0$, and this equals the claimed inequality only when $x+y\ge0$.

For nonnegative $x$ and $y$ both conditions hold: $xy\ge0$, so $\sqrt{xy}$ is defined and the square root of $\dfrac{(x+y)^2}{4}\ge xy\ge0$ may be taken, and $x+y\ge0$, so $|x+y|=x+y$. Hence
\[\dfrac{x+y}{2}=\sqrt{\dfrac{(x+y)^2}{4}}\ge\sqrt{xy}\qquad(x,y\ge0).\]
Equality holds exactly when $\dfrac{(x+y)^2}{4}=xy$, that is, $(x-y)^2=0$, so when $x=y$. $\blacksquare$
